%! Piecewise Linear Learning Functions — Unit 8, Lesson 8.1 — Student Packet Cover Sheet
\documentclass[10pt]{article}
\usepackage{linalg-article}
\usepackage{linalg-boxes}
\usepackage{ltablex}
\keepXColumns

\begin{document}

% Full-bleed burgundy banner
\begin{tikzpicture}[remember picture, overlay]
  \fill[burgundy] (current page.north west)
    rectangle ([yshift=-0.9in]current page.north east);
\end{tikzpicture}
\vspace{-0.8in}

\begin{center}
  {\color{white}\textbf{\LARGE Linear Algebra}}\\[4pt]
  {\color{blushmid}\large\textbf{Unit 8 \quad Learning from Data}}\\[6pt]
  {\color{white}\textbf{\large Lesson 8.1 \quad Piecewise Linear Learning Functions}}
\end{center}

\vspace{0.22in}
\namedateperiod
\vspace{0.18in}

\begin{learningtargetbox}
\small
\begin{enumerate}[leftmargin=1.5em, itemsep=5pt]
  \item \ldots{} build a \textbf{piecewise-linear} function by \textbf{shifting} and \textbf{adding} ReLUs, and read off where each \textbf{fold} sits.
  \item \ldots{} find the slope on each straight piece (the running sum of the switched-on ReLUs) and count pieces: $N$ folds make $N+1$ pieces.
  \item \ldots{} explain how a layer $\mathrm{ReLU}(A\mathbf{v}+\mathbf{b})$ makes these functions --- the bias places each fold, the weights set each slope --- so a wider layer fits wigglier data.
\end{enumerate}
\end{learningtargetbox}

\vspace{0.14in}

\begin{tocbox}
\small
\renewcommand{\arraystretch}{1.5}
\begin{tabularx}{\linewidth}{@{} c l X r @{}}
\rowcolor{blush}
\textbf{\#} & \textbf{Component} & \textbf{Description} & \textbf{Score} \\
\midrule
1 & Warm-Up        & Spiral review: one ReLU $=$ one fold, shifting the fold, and adding ReLUs to change slope & \blank{1.2cm} \\
2 & Guided Notes   & Adding shifted ReLUs; the tent function; reading slopes; counting folds and pieces; a layer builds them & \blank{1.2cm} \\
3 & Group Activity & Shifted ReLUs, building tents, the fold-of-$ax+b$, and how many folds fit wigglier data & \blank{1.2cm} \\
4 & Exit Ticket    & Quick check: evaluate a piecewise-linear function, count folds \& pieces, and why folds beat a line & \blank{1.2cm} \\
5 & Homework       & Shifted ReLUs, building and reading piecewise-linear graphs, fitting data by loss, and a look ahead & \blank{1.2cm} \\
\midrule
  & \multicolumn{2}{r}{\textbf{Total}} & \blank{1.2cm} \\
\end{tabularx}
\end{tocbox}

\vspace{0.10in}

\begin{remindbox}
\small A single \textbf{ReLU} bends once, at one \textbf{fold}. Shift it --- $\mathrm{ReLU}(x-s)$ --- and the
fold slides to $x=s$; scale it --- $c\cdot\mathrm{ReLU}(x-s)$ --- and $c$ sets how much the slope changes
there. \textbf{Add} several together and you get a \textbf{piecewise-linear} function: straight pieces joined
at folds, with $N$ folds giving $N+1$ pieces. This is exactly what one \textbf{layer}
$\mathrm{ReLU}(A\mathbf{v}+\mathbf{b})$ builds --- the \textbf{bias} places each fold, the \textbf{weights}
set each slope --- so a wider layer bends more and fits wigglier data.
\end{remindbox}

\end{document}
